\chapter{英文答题句式与模板}
\label{ch:UappC-english}

高级计量经济学试卷为全英文命题，宜用英文作答。本附录按“陈述模型与假设 $\to$ 写推导 $\to$ 陈述结论 $\to$ 检验决策 $\to$ 估计量性质 $\to$ 方法比较 $\to$ 识别假设”七类给出可直接套用的规范英文句子，左栏为英文，右栏为中文提示；每类再按主题分小节，便于从目录直接查找。最后两节给出考场作答结构模板（Setup $\to$ Derivation $\to$ Result $\to$ Interpretation/Assumptions）、一道 2025 年真题的示范骨架，以及常用连接词和易错表达。正文各节末尾的 English version 盒是按专题展开的版本，可配合使用。

%======================================================================
\section{陈述模型与假设（Stating the model and assumptions）}\label{sec:UappC-model}

\subsection{截面线性回归：模型、矩阵形式与标准假设}

{\renewcommand{\arraystretch}{1.15}\begin{longtable}{@{}>{\raggedright\arraybackslash}p{11.4cm}>{\raggedright\arraybackslash\small}p{4.3cm}@{}}
\toprule
\textbf{英文句式（可直接套用，按需替换符号）} & \textbf{中文提示 / 用法}\\
\midrule
\endfirsthead
\toprule
\textbf{英文句式（可直接套用，按需替换符号）} & \textbf{中文提示 / 用法}\\
\midrule
\endhead
\bottomrule
\endlastfoot
Consider the linear regression model $y_i=x_i\T\beta+u_i$, $i=1,\dots,n$, where $x_i$ is a $K\times 1$ vector of regressors including a constant and $u_i$ is an unobserved error term. & 写出模型并交代维度、含义（答题的第一句）\\
\addlinespace[2pt]
In matrix notation, $Y=X\beta+u$, where $Y$ is $n\times 1$, $X$ is $n\times K$ with $i$-th row $x_i\T$, and $u$ is $n\times 1$. & 矩阵形式；先写维度可避免后面出错\\
\addlinespace[2pt]
Assume that $\{(y_i,x_i)\}_{i=1}^n$ are i.i.d., $\E(u_i\mid x_i)=0$, $\E(x_ix_i\T)$ is positive definite, and the fourth moments of $(y_i,x_i)$ are finite. & OLS 大样本理论的标准假设\\
\addlinespace[2pt]
We further assume conditional homoskedasticity, $\E(u_i^2\mid x_i)=\sigma^2$ for all $i$. & 同方差假设（Gauss–Markov 定理、2SLS 有效时需要）\\
\end{longtable}}


\subsection{内生性与工具变量}

{\renewcommand{\arraystretch}{1.15}\begin{longtable}{@{}>{\raggedright\arraybackslash}p{11.4cm}>{\raggedright\arraybackslash\small}p{4.3cm}@{}}
\toprule
\textbf{英文句式（可直接套用，按需替换符号）} & \textbf{中文提示 / 用法}\\
\midrule
\endfirsthead
\toprule
\textbf{英文句式（可直接套用，按需替换符号）} & \textbf{中文提示 / 用法}\\
\midrule
\endhead
\bottomrule
\endlastfoot
The regressor $x_{1i}$ is endogenous, i.e.\ $\E(x_{1i}u_i)\neq 0$, so OLS is inconsistent. & 声明内生性\\
\addlinespace[2pt]
Let $z_i$ be an $L\times 1$ vector of instruments ($L\ge K$) satisfying the exogeneity condition $\E(z_iu_i)=0$ and the relevance (rank) condition $\rank\E(z_ix_i\T)=K$. & IV 的两个条件\\
\end{longtable}}


\subsection{面板模型：静态面板、动态面板与因子结构}

{\renewcommand{\arraystretch}{1.15}\begin{longtable}{@{}>{\raggedright\arraybackslash}p{11.4cm}>{\raggedright\arraybackslash\small}p{4.3cm}@{}}
\toprule
\textbf{英文句式（可直接套用，按需替换符号）} & \textbf{中文提示 / 用法}\\
\midrule
\endfirsthead
\toprule
\textbf{英文句式（可直接套用，按需替换符号）} & \textbf{中文提示 / 用法}\\
\midrule
\endhead
\bottomrule
\endlastfoot
Consider the static panel data model $y_{it}=\alpha_i+x_{it}\T\beta+u_{it}$, $i=1,\dots,N$, $t=1,\dots,T$, where $\alpha_i$ is an unobserved individual effect that may be correlated with $x_{it}$. & FE 模型（$\alpha_i$ 可与 $x_{it}$ 相关）\\
\addlinespace[2pt]
We assume strict exogeneity, $\E(u_{it}\mid x_{i1},\dots,x_{iT},\alpha_i)=0$ for all $t$, and consider asymptotics with $N\to\infty$ and $T$ fixed. & 严格外生 + 交代渐近框架\\
\addlinespace[2pt]
In the dynamic panel model $y_{it}=\alpha_i+\rho\,y_{i,t-1}+u_{it}$ with $|\rho|<1$, the errors are serially uncorrelated and sequentially exogenous: $\E(u_{it}\mid y_{i,t-1},\dots,y_{i0},\alpha_i)=0$. & 动态面板的序贯外生\\
\addlinespace[2pt]
Suppose the error has a factor structure, $u_{it}=\lambda_i\T f_t+v_{it}$, where $f_t$ is an $r\times 1$ vector of unobserved common factors and $\lambda_i$ is the vector of factor loadings. & 因子结构 / 交互固定效应\\
\end{longtable}}


\subsection{潜在结果框架}

{\renewcommand{\arraystretch}{1.15}\begin{longtable}{@{}>{\raggedright\arraybackslash}p{11.4cm}>{\raggedright\arraybackslash\small}p{4.3cm}@{}}
\toprule
\textbf{英文句式（可直接套用，按需替换符号）} & \textbf{中文提示 / 用法}\\
\midrule
\endfirsthead
\toprule
\textbf{英文句式（可直接套用，按需替换符号）} & \textbf{中文提示 / 用法}\\
\midrule
\endhead
\bottomrule
\endlastfoot
Let $y_{it}(1)$ and $y_{it}(0)$ denote the potential outcomes with and without treatment, so that the observed outcome is $y_{it}=D_{it}\,y_{it}(1)+(1-D_{it})\,y_{it}(0)$. & 潜在结果框架\\
\end{longtable}}


%======================================================================
\section{写推导（Writing derivations）}\label{sec:UappC-deriv}

\subsection{开头：说明要证什么、思路是什么}

{\renewcommand{\arraystretch}{1.15}\begin{longtable}{@{}>{\raggedright\arraybackslash}p{11.4cm}>{\raggedright\arraybackslash\small}p{4.3cm}@{}}
\toprule
\textbf{英文句式（可直接套用，按需替换符号）} & \textbf{中文提示 / 用法}\\
\midrule
\endfirsthead
\toprule
\textbf{英文句式（可直接套用，按需替换符号）} & \textbf{中文提示 / 用法}\\
\midrule
\endhead
\bottomrule
\endlastfoot
To show that $\hat\beta$ is consistent and asymptotically normal, I first write the estimation error in terms of sample averages and then apply the LLN and the CLT. & 先说“要证什么、思路是什么”\\
\end{longtable}}


\subsection{代入模型、取条件期望与迭代期望律}

{\renewcommand{\arraystretch}{1.15}\begin{longtable}{@{}>{\raggedright\arraybackslash}p{11.4cm}>{\raggedright\arraybackslash\small}p{4.3cm}@{}}
\toprule
\textbf{英文句式（可直接套用，按需替换符号）} & \textbf{中文提示 / 用法}\\
\midrule
\endfirsthead
\toprule
\textbf{英文句式（可直接套用，按需替换符号）} & \textbf{中文提示 / 用法}\\
\midrule
\endhead
\bottomrule
\endlastfoot
Substituting $Y=X\beta+u$ into the estimator gives $\hat\beta=\beta+(X\T X)\inv X\T u$. & 代入模型\\
\addlinespace[2pt]
Taking expectations conditional on $X$ and using $\E(u\mid X)=0$, we obtain $\E(\hat\beta\mid X)=\beta$. & 取条件期望\\
\addlinespace[2pt]
By the law of iterated expectations, $\E(x_iu_i)=\E[x_i\,\E(u_i\mid x_i)]=0$. & 迭代期望律 + 提出已知因子\\
\end{longtable}}


\subsection{大样本推导：分母用 LLN，分子用 CLT}

{\renewcommand{\arraystretch}{1.15}\begin{longtable}{@{}>{\raggedright\arraybackslash}p{11.4cm}>{\raggedright\arraybackslash\small}p{4.3cm}@{}}
\toprule
\textbf{英文句式（可直接套用，按需替换符号）} & \textbf{中文提示 / 用法}\\
\midrule
\endfirsthead
\toprule
\textbf{英文句式（可直接套用，按需替换符号）} & \textbf{中文提示 / 用法}\\
\midrule
\endhead
\bottomrule
\endlastfoot
Multiplying by $\sqrt n$, $\sqrt n(\hat\beta-\beta)=\bigl(\frac1n\sum_i x_ix_i\T\bigr)\inv\frac{1}{\sqrt n}\sum_i x_iu_i$. & 乘 $\sqrt n$，分配 $\frac1n$ 与 $\frac{1}{\sqrt n}$\\
\addlinespace[2pt]
By the WLLN, $\frac1n\sum_i x_ix_i\T\pto Q_{xx}$; by the CLT, $\frac{1}{\sqrt n}\sum_i x_iu_i\dto N(0,\Omega)$; the result then follows from the continuous mapping theorem and Slutsky's theorem. & “分母用 LLN，分子用 CLT”\\
\end{longtable}}


\subsection{一阶条件、零化矩阵、标量转置与迹技巧}

{\renewcommand{\arraystretch}{1.15}\begin{longtable}{@{}>{\raggedright\arraybackslash}p{11.4cm}>{\raggedright\arraybackslash\small}p{4.3cm}@{}}
\toprule
\textbf{英文句式（可直接套用，按需替换符号）} & \textbf{中文提示 / 用法}\\
\midrule
\endfirsthead
\toprule
\textbf{英文句式（可直接套用，按需替换符号）} & \textbf{中文提示 / 用法}\\
\midrule
\endhead
\bottomrule
\endlastfoot
The first-order condition $\partial J(\beta)/\partial\beta=0$ is equivalent to $X\T ZWZ\T(Y-X\beta)=0$, which yields $\hat\beta=(X\T ZWZ\T X)\inv X\T ZWZ\T Y$. & 写一阶条件并解出\\
\addlinespace[2pt]
Premultiplying the stacked model by $M_D=I_N\otimes M_T$ eliminates $\alpha_i$, because $M_T\one_T=0$. & 左乘零化矩阵消去效应\\
\addlinespace[2pt]
Note that $Y\T X\beta$ is a scalar and therefore equals its transpose $\beta\T X\T Y$. & 标量转置\\
\addlinespace[2pt]
Since $u\T Mu$ is a scalar and $\E(uu\T\mid X)=\sigma^2I_n$, $\E(u\T Mu\mid X)=\E[\tr(Muu\T)\mid X]=\tr\bigl(M\,\E(uu\T\mid X)\bigr)=\sigma^2\tr(M)=\sigma^2(n-K)$. & 迹技巧（$\tr(M)=n-K$）\\
\end{longtable}}


\subsection{“只需证明”与“不失一般性”}

{\renewcommand{\arraystretch}{1.15}\begin{longtable}{@{}>{\raggedright\arraybackslash}p{11.4cm}>{\raggedright\arraybackslash\small}p{4.3cm}@{}}
\toprule
\textbf{英文句式（可直接套用，按需替换符号）} & \textbf{中文提示 / 用法}\\
\midrule
\endfirsthead
\toprule
\textbf{英文句式（可直接套用，按需替换符号）} & \textbf{中文提示 / 用法}\\
\midrule
\endhead
\bottomrule
\endlastfoot
It suffices to show that the remaining terms are $o_p(1)$; without loss of generality, we normalize $F\T F/T=I_r$. & “只需证明”“不失一般性”\\
\end{longtable}}


%======================================================================
\section{陈述定理与结论（Stating results）}\label{sec:UappC-results}

\subsection{无偏、一致与渐近正态}

{\renewcommand{\arraystretch}{1.15}\begin{longtable}{@{}>{\raggedright\arraybackslash}p{11.4cm}>{\raggedright\arraybackslash\small}p{4.3cm}@{}}
\toprule
\textbf{英文句式（可直接套用，按需替换符号）} & \textbf{中文提示 / 用法}\\
\midrule
\endfirsthead
\toprule
\textbf{英文句式（可直接套用，按需替换符号）} & \textbf{中文提示 / 用法}\\
\midrule
\endhead
\bottomrule
\endlastfoot
Hence $\hat\beta$ is conditionally unbiased: $\E(\hat\beta\mid X)=\beta$. & 无偏\\
\addlinespace[2pt]
Therefore $\hat\beta\pto\beta$ as $n\to\infty$; that is, $\hat\beta$ is consistent. & 一致\\
\addlinespace[2pt]
$\sqrt n(\hat\beta-\beta)\dto N(0,V)$ with $V=Q_{xx}\inv\Omega Q_{xx}\inv$, which is consistently estimated by $\hat V=\hat Q_{xx}\inv\hat\Omega\hat Q_{xx}\inv$. & 渐近正态 + 方差估计\\
\end{longtable}}


\subsection{有效性：同方差下的夹心塌缩、Gauss–Markov 与有效 GMM}

{\renewcommand{\arraystretch}{1.15}\begin{longtable}{@{}>{\raggedright\arraybackslash}p{11.4cm}>{\raggedright\arraybackslash\small}p{4.3cm}@{}}
\toprule
\textbf{英文句式（可直接套用，按需替换符号）} & \textbf{中文提示 / 用法}\\
\midrule
\endfirsthead
\toprule
\textbf{英文句式（可直接套用，按需替换符号）} & \textbf{中文提示 / 用法}\\
\midrule
\endhead
\bottomrule
\endlastfoot
Under homoskedasticity the sandwich collapses to $V^0=\sigma^2Q_{xx}\inv$. & 夹心塌缩\\
\addlinespace[2pt]
By the Gauss--Markov theorem, OLS is the best linear unbiased estimator (BLUE) under the Gauss--Markov assumptions, in particular $\E(u\mid X)=0$ and $\Var(u\mid X)=\sigma^2I_n$. & BLUE\\
\addlinespace[2pt]
The efficient GMM estimator, which uses $W=\Omega\inv$, attains the asymptotic variance bound $(Q\T\Omega\inv Q)\inv$, where $Q=\E(z_ix_i\T)$ and $\Omega=\E(z_iz_i\T u_i^2)$. & 有效 GMM（$Q$ 为 $L\times K$，$\Omega$ 为 $L\times L$）\\
\end{longtable}}


\subsection{面板结论：数值等价、RE 收敛于 FE、Nickell 偏误与伴随参数问题}

{\renewcommand{\arraystretch}{1.15}\begin{longtable}{@{}>{\raggedright\arraybackslash}p{11.4cm}>{\raggedright\arraybackslash\small}p{4.3cm}@{}}
\toprule
\textbf{英文句式（可直接套用，按需替换符号）} & \textbf{中文提示 / 用法}\\
\midrule
\endfirsthead
\toprule
\textbf{英文句式（可直接套用，按需替换符号）} & \textbf{中文提示 / 用法}\\
\midrule
\endhead
\bottomrule
\endlastfoot
This shows that $\hat\beta_{LSDV}$ and $\hat\beta_{WG}$ are numerically identical. & 数值等价（不是“渐近等价”）\\
\addlinespace[2pt]
As $T\to\infty$, $\psi=\sigma_u^2/(\sigma_u^2+T\sigma_\alpha^2)\to 0$ and the RE (GLS) estimator converges to the within estimator. & RE $\to$ FE\\
\addlinespace[2pt]
For fixed $T$ the within estimator of $\rho$ is inconsistent, with asymptotic bias approximately $-(1+\rho)/(T-1)$ (the Nickell bias). & Nickell 偏误\\
\addlinespace[2pt]
$\hat\beta$ is consistent as $N\to\infty$, whereas the estimated individual effects $\hat\alpha_i$ are consistent only as $T\to\infty$ (the incidental parameter problem). & 伴随参数问题\\
\end{longtable}}


%======================================================================
\section{做检验决策（Making test decisions）}\label{sec:UappC-tests}

\subsection{写出假设与检验统计量}

{\renewcommand{\arraystretch}{1.15}\begin{longtable}{@{}>{\raggedright\arraybackslash}p{11.4cm}>{\raggedright\arraybackslash\small}p{4.3cm}@{}}
\toprule
\textbf{英文句式（可直接套用，按需替换符号）} & \textbf{中文提示 / 用法}\\
\midrule
\endfirsthead
\toprule
\textbf{英文句式（可直接套用，按需替换符号）} & \textbf{中文提示 / 用法}\\
\midrule
\endhead
\bottomrule
\endlastfoot
We test $H_0:R\beta=r$ against $H_1:R\beta\neq r$, where $R$ is a $q\times K$ matrix of rank $q$. & 先写假设与 $R$ 的维度\\
\addlinespace[2pt]
Under $H_0$, $W=(R\hat\beta-r)\T\bigl[R\hat V_{\hat\beta}R\T\bigr]\inv(R\hat\beta-r)\dto\chi^2_q$. & 写出统计量及其分布\\
\end{longtable}}


\subsection{拒绝、不拒绝与置信区间的标准说法}

{\renewcommand{\arraystretch}{1.15}\begin{longtable}{@{}>{\raggedright\arraybackslash}p{11.4cm}>{\raggedright\arraybackslash\small}p{4.3cm}@{}}
\toprule
\textbf{英文句式（可直接套用，按需替换符号）} & \textbf{中文提示 / 用法}\\
\midrule
\endfirsthead
\toprule
\textbf{英文句式（可直接套用，按需替换符号）} & \textbf{中文提示 / 用法}\\
\midrule
\endhead
\bottomrule
\endlastfoot
Since $W=9.10>\chi^2_{3,0.95}=7.815$, we reject $H_0$ at the 5\% significance level ($p=0.028$). & 拒绝的标准说法（临界值查附录 A）\\
\addlinespace[2pt]
Because $|t|=1.42<1.96$, we fail to reject $H_0$ at the 5\% level; this does not prove that $H_0$ is true. & 用 “fail to reject”，不要说 “accept”\\
\addlinespace[2pt]
The 95\% confidence interval $\hat\beta_j\pm 1.96\,\mathrm{se}(\hat\beta_j)$ excludes zero, so $\beta_j$ is statistically significant at the 5\% level. & 用置信区间作判断\\
\end{longtable}}


\subsection{常见检验结果的解读：弱工具、\texorpdfstring{$J$}{J} 检验、Hausman、AR(2) 与 CD}

{\renewcommand{\arraystretch}{1.15}\begin{longtable}{@{}>{\raggedright\arraybackslash}p{11.4cm}>{\raggedright\arraybackslash\small}p{4.3cm}@{}}
\toprule
\textbf{英文句式（可直接套用，按需替换符号）} & \textbf{中文提示 / 用法}\\
\midrule
\endfirsthead
\toprule
\textbf{英文句式（可直接套用，按需替换符号）} & \textbf{中文提示 / 用法}\\
\midrule
\endhead
\bottomrule
\endlastfoot
The first-stage $F$ statistic (35.2) exceeds the rule-of-thumb threshold of 10 and the Stock--Yogo critical value, so the instruments are not weak. & 相关性检验\\
\addlinespace[2pt]
The Hansen $J$ statistic is 1.80 with $L-K=2$ degrees of freedom ($p=0.41$), so we do not reject the overidentifying restrictions: there is no evidence against instrument validity. & 过度识别检验（不拒绝 $\neq$ 工具有效）\\
\addlinespace[2pt]
The Hausman statistic rejects the null that $\alpha_i$ is uncorrelated with $x_{it}$; hence the RE estimator is inconsistent and we use the FE estimator. & FE vs RE\\
\addlinespace[2pt]
The AR(1) test rejects while the AR(2) test does not, which is consistent with serially uncorrelated $u_{it}$ and supports the validity of $y_{i,t-2}$ as an instrument. & 动态面板设定检验\\
\addlinespace[2pt]
The Pesaran CD statistic (12.3) is far above 1.96, indicating strong cross-sectional dependence, so a model with common factors is needed. & 截面相关检验\\
\end{longtable}}


%======================================================================
\section{解释估计量性质（Explaining properties of estimators）}\label{sec:UappC-props}

\subsection{无偏与一致；异方差与聚类标准误}

{\renewcommand{\arraystretch}{1.15}\begin{longtable}{@{}>{\raggedright\arraybackslash}p{11.4cm}>{\raggedright\arraybackslash\small}p{4.3cm}@{}}
\toprule
\textbf{英文句式（可直接套用，按需替换符号）} & \textbf{中文提示 / 用法}\\
\midrule
\endfirsthead
\toprule
\textbf{英文句式（可直接套用，按需替换符号）} & \textbf{中文提示 / 用法}\\
\midrule
\endhead
\bottomrule
\endlastfoot
Unbiasedness is a finite-sample property, whereas consistency is a large-sample property; neither implies the other. & 无偏与一致的区别\\
\addlinespace[2pt]
Under heteroskedasticity OLS remains unbiased and consistent but is no longer efficient, and the conventional standard errors are invalid; we therefore report heteroskedasticity-robust (White) standard errors. & 异方差的后果与对策\\
\addlinespace[2pt]
Cluster-robust standard errors allow arbitrary correlation within clusters but require the number of clusters $G$ to be large; the point estimates are unchanged. & 聚类标准误\\
\end{longtable}}


\subsection{2SLS 的有限样本偏误与弱工具}

{\renewcommand{\arraystretch}{1.15}\begin{longtable}{@{}>{\raggedright\arraybackslash}p{11.4cm}>{\raggedright\arraybackslash\small}p{4.3cm}@{}}
\toprule
\textbf{英文句式（可直接套用，按需替换符号）} & \textbf{中文提示 / 用法}\\
\midrule
\endfirsthead
\toprule
\textbf{英文句式（可直接套用，按需替换符号）} & \textbf{中文提示 / 用法}\\
\midrule
\endhead
\bottomrule
\endlastfoot
2SLS is consistent but biased in finite samples; with weak instruments it is biased toward OLS and its sampling distribution is far from normal. & 弱工具\\
\end{longtable}}


\subsection{FE、RE 与 Arellano–Bond 估计量的性质}

{\renewcommand{\arraystretch}{1.15}\begin{longtable}{@{}>{\raggedright\arraybackslash}p{11.4cm}>{\raggedright\arraybackslash\small}p{4.3cm}@{}}
\toprule
\textbf{英文句式（可直接套用，按需替换符号）} & \textbf{中文提示 / 用法}\\
\midrule
\endfirsthead
\toprule
\textbf{英文句式（可直接套用，按需替换符号）} & \textbf{中文提示 / 用法}\\
\midrule
\endhead
\bottomrule
\endlastfoot
The within estimator uses only within-unit variation over time: it is robust to any correlation between $\alpha_i$ and $x_{it}$, but it cannot identify the coefficients of time-invariant regressors. & FE 的优缺点\\
\addlinespace[2pt]
The RE estimator is efficient under the RE assumptions but inconsistent if $\E(\alpha_i\mid x_{it})\neq 0$. & RE 的优缺点\\
\addlinespace[2pt]
The Arellano--Bond estimator is consistent for fixed $T$ as $N\to\infty$, but lagged levels become weak instruments when $\rho$ is close to one. & AB 的性质\\
\end{longtable}}


\subsection{因子的旋转不确定性与个体处理效应估计}

{\renewcommand{\arraystretch}{1.15}\begin{longtable}{@{}>{\raggedright\arraybackslash}p{11.4cm}>{\raggedright\arraybackslash\small}p{4.3cm}@{}}
\toprule
\textbf{英文句式（可直接套用，按需替换符号）} & \textbf{中文提示 / 用法}\\
\midrule
\endfirsthead
\toprule
\textbf{英文句式（可直接套用，按需替换符号）} & \textbf{中文提示 / 用法}\\
\midrule
\endhead
\bottomrule
\endlastfoot
The factors and loadings are identified only up to an invertible rotation, but the common component $\lambda_i\T f_t$ and $\hat\beta$ are invariant to the rotation. & 因子识别\\
\addlinespace[2pt]
The estimator of an individual treatment effect is unbiased but not consistent, because its error contains the idiosyncratic shock of that unit and period. & 个体处理效应\\
\end{longtable}}


%======================================================================
\section{比较方法（Comparing methods）}\label{sec:UappC-compare}

\subsection{OLS 与 IV；2SLS 与两步 GMM}

{\renewcommand{\arraystretch}{1.15}\begin{longtable}{@{}>{\raggedright\arraybackslash}p{11.4cm}>{\raggedright\arraybackslash\small}p{4.3cm}@{}}
\toprule
\textbf{英文句式（可直接套用，按需替换符号）} & \textbf{中文提示 / 用法}\\
\midrule
\endfirsthead
\toprule
\textbf{英文句式（可直接套用，按需替换符号）} & \textbf{中文提示 / 用法}\\
\midrule
\endhead
\bottomrule
\endlastfoot
Compared with OLS, IV is consistent under endogeneity but less precise: with a single instrument its asymptotic variance is inflated by the factor $1/\rho_{zx}^2$. & OLS vs IV\\
\addlinespace[2pt]
2SLS is GMM with $W=(Z\T Z)\inv$; it is efficient under conditional homoskedasticity, whereas two-step GMM is efficient under general heteroskedasticity. & 2SLS vs 两步 GMM\\
\end{longtable}}


\subsection{FE 与 RE；差分 GMM 与系统 GMM}

{\renewcommand{\arraystretch}{1.15}\begin{longtable}{@{}>{\raggedright\arraybackslash}p{11.4cm}>{\raggedright\arraybackslash\small}p{4.3cm}@{}}
\toprule
\textbf{英文句式（可直接套用，按需替换符号）} & \textbf{中文提示 / 用法}\\
\midrule
\endfirsthead
\toprule
\textbf{英文句式（可直接套用，按需替换符号）} & \textbf{中文提示 / 用法}\\
\midrule
\endhead
\bottomrule
\endlastfoot
FE is consistent under both the FE and the RE models, whereas RE is consistent only under the RE model but is more efficient when that model is correct. & FE vs RE\\
\addlinespace[2pt]
Difference GMM uses lagged levels as instruments for the differenced equation, while System GMM additionally uses lagged differences as instruments for the level equation, at the cost of an initial-condition assumption. & AB vs System\\
\end{longtable}}


\subsection{Bai 迭代主成分与 CCE}

{\renewcommand{\arraystretch}{1.15}\begin{longtable}{@{}>{\raggedright\arraybackslash}p{11.4cm}>{\raggedright\arraybackslash\small}p{4.3cm}@{}}
\toprule
\textbf{英文句式（可直接套用，按需替换符号）} & \textbf{中文提示 / 用法}\\
\midrule
\endfirsthead
\toprule
\textbf{英文句式（可直接套用，按需替换符号）} & \textbf{中文提示 / 用法}\\
\midrule
\endhead
\bottomrule
\endlastfoot
Bai's (2009) iterative principal components estimator requires both $N$ and $T$ to be large and the number of factors to be known, whereas the CCE estimator only needs $N\to\infty$ and does not require $r$, but assumes that the regressors load linearly on the factors. & Bai vs CCE\\
\end{longtable}}


\subsection{SCM 与 PDA；平均效应与个体效应；交错处理}

{\renewcommand{\arraystretch}{1.15}\begin{longtable}{@{}>{\raggedright\arraybackslash}p{11.4cm}>{\raggedright\arraybackslash\small}p{4.3cm}@{}}
\toprule
\textbf{英文句式（可直接套用，按需替换符号）} & \textbf{中文提示 / 用法}\\
\midrule
\endfirsthead
\toprule
\textbf{英文句式（可直接套用，按需替换符号）} & \textbf{中文提示 / 用法}\\
\midrule
\endhead
\bottomrule
\endlastfoot
SCM restricts the weights to be non-negative and to sum to one (a convex combination of control units), whereas PDA uses an unrestricted linear combination estimated by OLS on the pre-treatment periods. & SCM vs PDA\\
\addlinespace[2pt]
DiD and TWFE target the average treatment effect on the treated, whereas SCM and PDA estimate the unit-specific effects of a single (or a few) treated unit(s). & 平均效应 vs 个体效应\\
\addlinespace[2pt]
With staggered adoption and heterogeneous effects, static TWFE may put negative weights on some cohort-time effects; heterogeneity-robust estimators (Callaway--Sant'Anna, Sun--Abraham, imputation) use only clean comparisons with not-yet- or never-treated units. & 交错处理\\
\end{longtable}}


%======================================================================
\section{讨论识别假设（Discussing identifying assumptions）}\label{sec:UappC-ident}

\subsection{IV：相关性、外生性与恰好识别时不可检验}

{\renewcommand{\arraystretch}{1.15}\begin{longtable}{@{}>{\raggedright\arraybackslash}p{11.4cm}>{\raggedright\arraybackslash\small}p{4.3cm}@{}}
\toprule
\textbf{英文句式（可直接套用，按需替换符号）} & \textbf{中文提示 / 用法}\\
\midrule
\endfirsthead
\toprule
\textbf{英文句式（可直接套用，按需替换符号）} & \textbf{中文提示 / 用法}\\
\midrule
\endhead
\bottomrule
\endlastfoot
Identification requires the instruments to be relevant, $\rank\E(z_ix_i\T)=K$, and exogenous, $\E(z_iu_i)=0$; the exclusion restriction cannot be tested when the model is exactly identified. & IV 识别；恰好识别时无法检验\\
\end{longtable}}


\subsection{DiD：平行趋势、无预期与 SUTVA、事前检验的局限}

{\renewcommand{\arraystretch}{1.15}\begin{longtable}{@{}>{\raggedright\arraybackslash}p{11.4cm}>{\raggedright\arraybackslash\small}p{4.3cm}@{}}
\toprule
\textbf{英文句式（可直接套用，按需替换符号）} & \textbf{中文提示 / 用法}\\
\midrule
\endfirsthead
\toprule
\textbf{英文句式（可直接套用，按需替换符号）} & \textbf{中文提示 / 用法}\\
\midrule
\endhead
\bottomrule
\endlastfoot
The key identifying assumption is parallel trends: in the absence of treatment, the average outcomes of the treated and control groups would have followed the same trend. & 平行趋势\\
\addlinespace[2pt]
We also assume no anticipation and SUTVA (no spillovers between units); spillovers onto the control group would bias the DiD estimate. & 无预期、SUTVA\\
\addlinespace[2pt]
Jointly insignificant pre-treatment coefficients support, but do not prove, the parallel trends assumption. & 事前检验的局限\\
\end{longtable}}


\subsection{面板：FE 的识别来源与系统 GMM 的初始条件}

{\renewcommand{\arraystretch}{1.15}\begin{longtable}{@{}>{\raggedright\arraybackslash}p{11.4cm}>{\raggedright\arraybackslash\small}p{4.3cm}@{}}
\toprule
\textbf{英文句式（可直接套用，按需替换符号）} & \textbf{中文提示 / 用法}\\
\midrule
\endfirsthead
\toprule
\textbf{英文句式（可直接套用，按需替换符号）} & \textbf{中文提示 / 用法}\\
\midrule
\endhead
\bottomrule
\endlastfoot
The FE estimator is identified from within-unit variation over time; regressors that do not vary over time are absorbed by $\alpha_i$. & FE 的识别来源\\
\addlinespace[2pt]
System GMM requires the mean-stationarity condition $y_{i0}=\alpha_i/(1-\rho)+v_{i0}$ with $v_{i0}$ uncorrelated with $\alpha_i$, so that $\Delta y_{i,t-1}$ is uncorrelated with $\alpha_i$. & System GMM 的额外假设\\
\end{longtable}}


\subsection{因子标准化、SCM 的适用条件与 FEF 的识别}

{\renewcommand{\arraystretch}{1.15}\begin{longtable}{@{}>{\raggedright\arraybackslash}p{11.4cm}>{\raggedright\arraybackslash\small}p{4.3cm}@{}}
\toprule
\textbf{英文句式（可直接套用，按需替换符号）} & \textbf{中文提示 / 用法}\\
\midrule
\endfirsthead
\toprule
\textbf{英文句式（可直接套用，按需替换符号）} & \textbf{中文提示 / 用法}\\
\midrule
\endhead
\bottomrule
\endlastfoot
Factors and loadings are identified only up to rotation, so we impose $r^2$ normalizations: $F\T F/T=I_r$ and $\Lambda\T\Lambda$ diagonal. & 因子标准化（$\frac{r(r+1)}{2}+\frac{r(r-1)}{2}=r^2$ 个约束）\\
\addlinespace[2pt]
SCM is valid if the loadings of the treated unit lie in the convex hull of the control units' loadings and the pre-treatment fit is good over a long pre-treatment period. & SCM 的条件\\
\addlinespace[2pt]
The FEF estimator identifies $\gamma$ only if the time-invariant regressors $z_i$ are uncorrelated with $\alpha_i$; otherwise instruments are needed in the second step (FEF-IV). & 时不变变量的识别\\
\end{longtable}}


%======================================================================
\section{考场作答结构模板}\label{sec:UappC-structure}

\subsection{四步结构：Setup \texorpdfstring{$\to$}{->} Derivation \texorpdfstring{$\to$}{->} Result \texorpdfstring{$\to$}{->} Interpretation/Assumptions}

{\renewcommand{\arraystretch}{1.15}\begin{longtable}{@{}>{\raggedright\arraybackslash}p{2.6cm}>{\raggedright\arraybackslash}p{6.2cm}>{\raggedright\arraybackslash}p{6.9cm}@{}}
\toprule
\textbf{步骤} & \textbf{要写什么（中文提示）} & \textbf{英文骨架}\\
\midrule
\endfirsthead
\toprule
\textbf{步骤} & \textbf{要写什么（中文提示）} & \textbf{英文骨架}\\
\midrule
\endhead
\bottomrule
\endlastfoot
1. Setup & 写模型、下标范围、维度；定义要用的记号（均值、变换矩阵等）；写出假设与渐近框架 & “Consider \ldots, where \ldots\ is $n\times K$. Assume \ldots. Define \ldots.”\\
\addlinespace[2pt]
2. Derivation & 先一句“要证什么、思路是什么”；按 hint 的路线逐步推导，每步写一个理由（代入、LIE、$M_T\one_T=0$、LLN/CLT 等） & “To show \ldots, I first \ldots. Substituting \ldots\ gives \ldots. By the LIE / WLLN / CLT, \ldots.”\\
\addlinespace[2pt]
3. Result & 把答案单独写成一行（方框或 “Hence”），给出估计量公式 / 分布 / 等价关系 & “Hence $\hat\beta=\ldots$.” “Therefore $\sqrt n(\hat\beta-\beta)\dto N(0,V)$.”\\
\addlinespace[2pt]
4. Interpretation / Assumptions & 解释结论的含义；列出所需假设（外生性、秩条件、平衡面板等）；指出假设失效时的后果与对策 & “This shows that \ldots. The result requires \ldots. If \ldots\ fails, \ldots\ is inconsistent and we should use \ldots\ instead.”\\
\end{longtable}}


\begin{note}[作答要点]
(1) 题目给了 hint，就\key{顺着 hint 的路线写}；(2) 先定义后使用：凡与题目记号不同的符号，第一次出现时写一句 “where \ldots\ denotes \ldots”；(3) 推导不跳步，矩阵写维度；(4) 结论句与推导分开，让阅卷人一眼找到；(5) 时间不够时，至少写出 Setup、关键公式和结论，并说明所需假设。
\end{note}

\subsection{示范骨架：2025 年真题第 1(a) 题（双向组内变换）}

\tagexam\ 题目：对平衡面板 TWFE 模型 $Y_{it}=\alpha_i+\lambda_t+\beta x_{it}+\epsilon_{it}$，写出同时消去 $\alpha_i$ 与 $\lambda_t$ 的组内变换及 $\hat\beta_{FE}$（完整讲解见第 17 章）。

\begin{english}[Model answer skeleton]
\textbf{Setup.} Consider $Y_{it}=\alpha_i+\lambda_t+\beta x_{it}+\epsilon_{it}$, $i=1,\dots,N$, $t=1,\dots,T$, with a balanced panel. Define $\bar Y_{i\cdot}=\frac1T\sum_t Y_{it}$, $\bar Y_{\cdot t}=\frac1N\sum_i Y_{it}$, $\bar Y_{\cdot\cdot}=\frac{1}{NT}\sum_i\sum_t Y_{it}$, and similarly for $x$ and $\epsilon$; let $\bar\alpha=\frac1N\sum_i\alpha_i$ and $\bar\lambda=\frac1T\sum_t\lambda_t$.

\textbf{Derivation.} Averaging the model over $t$, over $i$, and over both gives
\begin{gather*}
\bar Y_{i\cdot}=\alpha_i+\bar\lambda+\beta\bar x_{i\cdot}+\bar\epsilon_{i\cdot},\qquad
\bar Y_{\cdot t}=\bar\alpha+\lambda_t+\beta\bar x_{\cdot t}+\bar\epsilon_{\cdot t},\\
\bar Y_{\cdot\cdot}=\bar\alpha+\bar\lambda+\beta\bar x_{\cdot\cdot}+\bar\epsilon_{\cdot\cdot}.
\end{gather*}
Hence the two-way within transformation $\ddot Y_{it}=Y_{it}-\bar Y_{i\cdot}-\bar Y_{\cdot t}+\bar Y_{\cdot\cdot}$ satisfies $\ddot Y_{it}=\beta\ddot x_{it}+\ddot\epsilon_{it}$, because $\alpha_i-\alpha_i-\bar\alpha+\bar\alpha=0$ and $\lambda_t-\bar\lambda-\lambda_t+\bar\lambda=0$.

\textbf{Result.} Applying OLS to the transformed equation,
\[
\hat\beta_{FE}=\frac{\sum_{i=1}^N\sum_{t=1}^T\ddot x_{it}\ddot Y_{it}}{\sum_{i=1}^N\sum_{t=1}^T\ddot x_{it}^2};
\qquad\text{in matrix form}\quad \hat\beta_{FE}=(x\T Qx)\inv x\T QY,\quad Q=M_N\otimes M_T ,
\]
where $x$ and $Y$ are the $NT\times 1$ vectors stacking $x_{it}$ and $Y_{it}$ by individual (with $t$ running fastest), and $M_m=I_m-m^{-1}\one_m\one_m\T$ is the $m\times m$ demeaning matrix.

\textbf{Interpretation and assumptions.} Both $\alpha_i$ and $\lambda_t$ are removed exactly because the panel is balanced. $\hat\beta_{FE}$ is unbiased under strict exogeneity, $\E(\epsilon_{it}\mid x_{i1},\dots,x_{iT},\alpha_i,\lambda_t)=0$, together with independence across $i$, and it is consistent as $N\to\infty$ with $T$ fixed, provided that $\frac{1}{NT}\sum_i\sum_t\ddot x_{it}^2$ has a positive probability limit. It is numerically identical to the estimator obtained from the one-way within transformation with time dummies (FWL theorem). In an unbalanced panel the simple double-demeaning formula no longer removes $\alpha_i+\lambda_t$; one should then include unit and time dummies (equivalently, use the Wansbeek--Kapteyn two-way transformation or iterated demeaning).
\end{english}

%======================================================================
\section{常用连接词与易错表达}\label{sec:UappC-phrases}

\subsection{推导与论证中的连接词}

{\renewcommand{\arraystretch}{1.15}\begin{longtable}{@{}>{\raggedright\arraybackslash}p{9.6cm}>{\raggedright\arraybackslash\small}p{6.1cm}@{}}
\toprule
\textbf{英文} & \textbf{用法}\\
\midrule
\endfirsthead
\toprule
\textbf{英文} & \textbf{用法}\\
\midrule
\endhead
\bottomrule
\endlastfoot
Hence / Therefore / Thus / It follows that & 引出结论\\
\addlinespace[2pt]
Since \ldots\ / Because \ldots\ / Given that \ldots & 给出理由（放在句首）\\
\addlinespace[2pt]
where \ldots\ denotes \ldots & 公式后解释记号\\
\addlinespace[2pt]
which implies that \ldots\ / so that \ldots & 推出下一步\\
\addlinespace[2pt]
under $H_0$ / under the null hypothesis & 检验分布的前提\\
\addlinespace[2pt]
as $N\to\infty$ with $T$ fixed / as $(N,T)\to\infty$ jointly & 交代渐近框架\\
\addlinespace[2pt]
provided that \ldots\ / as long as \ldots & 所需条件\\
\addlinespace[2pt]
in contrast / whereas / on the other hand & 比较方法\\
\addlinespace[2pt]
up to an invertible rotation / up to scale & “在……意义下”唯一\\
\addlinespace[2pt]
numerically identical / asymptotically equivalent & 数值相同 vs 渐近等价（不要混用）\\
\end{longtable}}


\subsection{常见错误表达与正确说法}

{\renewcommand{\arraystretch}{1.15}\begin{longtable}{@{}>{\raggedright\arraybackslash}p{4.6cm}>{\raggedright\arraybackslash}p{6.6cm}>{\raggedright\arraybackslash}p{4.5cm}@{}}
\toprule
\textbf{不要写} & \textbf{应写} & \textbf{理由}\\
\midrule
\endfirsthead
\toprule
\textbf{不要写} & \textbf{应写} & \textbf{理由}\\
\midrule
\endhead
\bottomrule
\endlastfoot
We accept $H_0$. & We fail to reject (do not reject) $H_0$. & 不拒绝不等于证明为真\\
\addlinespace[2pt]
The $p$-value is the probability that $H_0$ is true. & The $p$-value is the probability, under $H_0$, of obtaining a test statistic at least as extreme as the one observed. & $p$ 值的定义\\
\addlinespace[2pt]
The estimate is consistent. & The estimator is consistent. & 一致性是估计量（规则）的性质\\
\addlinespace[2pt]
$\hat\beta$ converges to $\beta$. & $\hat\beta$ converges in probability to $\beta$ ($\hat\beta\pto\beta$). & 写清收敛方式\\
\addlinespace[2pt]
The instruments are valid because $J$ is insignificant. & There is no evidence against instrument validity. & $J$ 检验只检验工具之间的相对有效性\\
\addlinespace[2pt]
Parallel trends hold because the pre-trend test is insignificant. & The pre-trend test supports, but does not prove, parallel trends. & 事前检验的功效可能很低\\
\addlinespace[2pt]
OLS is biased because of heteroskedasticity. & Under heteroskedasticity OLS is unbiased but inefficient, and the usual standard errors are wrong. & 异方差不影响无偏性\\
\addlinespace[2pt]
SSE（不加说明） & the residual sum of squares $SSR=\sum_i\hat u_i^2$ & 不同教材 SSE 的含义相反（Hansen 指残差平方和，Wooldridge 指解释平方和），第一次出现要写全称\\
\end{longtable}}
